% Generated draft: requires mathematical and source review.
\documentclass[11pt]{article}
\usepackage[margin=25mm]{geometry}
\usepackage{fontspec}
\setmainfont{Latin Modern Roman}
\usepackage{amsmath,amssymb,mathtools,unicode-math}
\setmathfont{Latin Modern Math}
\usepackage{enumitem,needspace,xurl,hyperref}
\setlength{\parindent}{0pt}
\setlength{\parskip}{5pt}
\setlength{\emergencystretch}{3em}
\newcommand{\N}{\mathbb{N}}
\newcommand{\Z}{\mathbb{Z}}
\newcommand{\Q}{\mathbb{Q}}
\newcommand{\R}{\mathbb{R}}
\newcommand{\C}{\mathbb{C}}
\newcommand{\F}{\mathbb{F}}
\newcommand{\A}{\mathbb{A}}
\newcommand{\PP}{\mathbb{P}}
\newcommand{\E}{\mathbb{E}}
\newcommand{\eps}{\varepsilon}
\newcommand{\dd}{\,\mathrm{d}}
\begin{document}
\section*{UnsolvedMath ID: 9400171}
\textbf{Source problem code:} catalogue entry
\subsection{Definitions and mathematical statement}
Fix an integer \(a>0\). Define the Lucas sequence for parameters \((a,-1)\) by
\[
 U_0=0,\qquad U_1=1,\qquad U_{n+2}=aU_{n+1}+U_n\quad(n\ge0),
\]
and put \(D=a^2+4\). Under the standard unramified convention, an odd prime \(p\nmid D\) is a Lucas--Wieferich prime for this pair if
\[
 U_{p-\left(\frac Dp\right)}\equiv0\pmod{p^2},
\]
where \(\left(\frac Dp\right)\in\{1,-1\}\) is the Legendre symbol. The usual Lucas congruence already gives divisibility by \(p\); the displayed condition asks for divisibility by \(p^2\).

The question is whether, for every fixed \(a>0\), there are infinitely many such primes. Explicitly,
\[
 \forall a\in\mathbb Z_{>0}\ \forall X>0\ \exists p>X\text{ prime}:\quad
 p\nmid2(a^2+4),\qquad
 p^2\mid U_{p-\left(\frac{a^2+4}{p}\right)}.
\]
When \(a=1\), \(U_n\) is the Fibonacci sequence; when \(a=2\), it is the Pell sequence. These give the two special cases identified by the source. The parameter \(a\) is fixed before the infinitude assertion; infinitude merely after varying \(a\) is weaker and is not the target.

The excerpt does not define the treatment of the finitely many primes dividing \(2D\). The displayed standard exclusion cannot affect infinitude for a fixed \(a\). The supplied summary labels the assertion open and says only finitely many examples are known, without a list or count; no positive number of known examples in either special case is inferred. The pinned Wikipedia label and later assessment await primary-source/current-status review.
\subsection{Short English statement}
For each positive integer \(a\), are there infinitely many primes where the Lucas congruence for \(U_{n+2}=aU_{n+1}+U_n\) strengthens from divisibility by \(p\) to \(p^2\)? Fibonacci and Pell Wieferich primes are the cases \(a=1\) and \(a=2\).
\subsection{Sources}
\begin{itemize}
\item Local registry: \url{data/unsolvedmath/problems.json}, permanent UnsolvedMath ID \(9400171\), source code catalogue entry; supplied statement, background, and research summary, updated 2026-08-06.
\item Wikipedia contributors, \emph{List of unsolved problems in mathematics}, Number theory, prime numbers bullet 171, wikitext line 205, pinned revision 1366636767: \url{https://en.wikipedia.org/w/index.php?title=List_of_unsolved_problems_in_mathematics&oldid=1366636767#Number_theory}. MediaWiki-source extraction accessed 2026-07-29; primary-source audit and final semantic deduplication remain pending.
\item The summary refers to the Wikipedia ``Lucas--Wieferich prime'' page without an original-paper citation or verified list of examples.
\end{itemize}
\end{document}
